% Diffusion Language Models are Faster than Autoregressive on CPU
% v2.1 — March 22, 2026
% Updated with Dream-7B comparative benchmarks and inter-step KV cache results
%
% NOTE: Figures 1–4 must be generated separately (see research/benchmark/
% scripts or regenerate with matplotlib). Place as fig1.pdf, fig2.pdf,
% fig3.pdf, fig4.pdf in the same directory, or comment out \includegraphics.

\documentclass[11pt]{article}

\usepackage[margin=1in]{geometry}
\usepackage{amsmath,amssymb}
\usepackage{booktabs}
\usepackage{graphicx}
\usepackage{hyperref}
\usepackage[numbers]{natbib}
\usepackage{xcolor}
% \usepackage{microtype}  % requires scalable fonts

\hypersetup{
  colorlinks=true,
  linkcolor=blue,
  citecolor=blue,
  urlcolor=blue
}

\title{Diffusion Language Models are Faster than Autoregressive on CPU:\\
An Empirical Study of the Memory-Compute Regime Inversion}

\author{Carmen Esteban\thanks{Corresponding author: \href{mailto:caresment@gmail.com}{caresment@gmail.com}. ORCID: 0009-0009-8062-5492}\\
IAFiscal \& Partners, Madrid, Spain}

\date{March 22, 2026}

\begin{document}

\maketitle

\begin{abstract}
Diffusion language models (dLLMs) are widely considered slower than autoregressive
(AR) models for inference. We demonstrate empirically that this conclusion is hardware-dependent:
on commodity CPUs, dLLMs can outperform optimized AR inference. We identify the
mechanism---a regime inversion from memory-bound (AR) to compute-bound (diffusion)---and
validate it with systematic benchmarks. We present \texttt{diffuse-cpp}, the first C++
inference engine for diffusion LLMs with GGML quantization support. On a 12-core AMD EPYC,
\texttt{diffuse-cpp} supports two models---LLaDA-8B (Llama backbone, MHA) and
Dream-7B (Qwen2.5 backbone, GQA)---and achieves 14--28~tok/s with entropy-exit
scheduling, inter-step KV cache, and 256-token generation on easy prompts, up to
$3.3\times$ faster than \texttt{llama.cpp} (8.51~tok/s) on the same hardware. Dream
excels at math and code (21.6~tok/s, correctly solving arithmetic), while LLaDA excels
at translation (27.7~tok/s). Thread scaling reaches $7.5\times$ from 1 to 12 cores for
diffusion versus saturation at 4 cores for AR. An entropy-based adaptive scheduler
reduces denoising steps from 16 to 2--4 on easy prompts without quality loss. An
inter-step KV cache reuses Key and Value tensors between denoising steps, providing
$1.6\times$ average speedup with implicit regularization benefits. These results suggest
that diffusion models may be a viable path toward democratizing LLM inference on commodity
hardware without GPU dependency.
\end{abstract}

\textbf{Keywords:} diffusion language models, CPU inference, memory-bound, compute-bound,
LLM acceleration, quantization, adaptive scheduling, inter-step cache, multi-model

\textbf{License:} CC-BY-4.0 (paper), Apache-2.0 (code)


\section{Introduction}

Large language model (LLM) inference is dominated by GPU hardware. The supply chain---from
ASML's lithography monopoly through TSMC's fabrication to NVIDIA's tensor cores---creates
a dependency that concentrates access to AI capabilities among well-funded actors. The
fundamental question motivating this work is whether alternative computational paradigms
can make LLM inference viable on commodity CPUs, which are ubiquitous and inexpensive.

Autoregressive (AR) models generate text one token at a time. Each token requires a full
pass through the model's weight matrix, reading billions of parameters from memory to
produce a single output. On CPUs, where memory bandwidth is $10$--$50\times$ lower than
on GPUs, this sequential pattern creates a severe bottleneck: the processor spends most
of its time waiting for data rather than computing. Adding more CPU cores does not help,
because all threads compete for the same memory bus. This is the memory-bound regime.

Diffusion language models (dLLMs)~\cite{nie2025,lou2024,sahoo2024} operate differently.
Instead of generating tokens sequentially, they begin with a fully masked sequence and
iteratively refine all positions in parallel across $T$ denoising steps. Each step reads
the model weights once but processes $B$ tokens simultaneously through matrix-matrix
multiplication rather than matrix-vector. This amortizes the cost of memory access across
many tokens, shifting the workload from memory-bound to compute-bound. Crucially,
matrix-matrix multiplication parallelizes efficiently across CPU cores.

Despite this favorable computational profile, the dLLM literature has focused exclusively
on GPU benchmarks. Papers introducing LLaDA~\cite{nie2025}, SEDD~\cite{lou2024}, and
MDLM~\cite{sahoo2024} report wall-clock times on A100 and H100 GPUs, where the massive
memory bandwidth makes AR models fast and the compute advantages of diffusion are less
pronounced. No prior work has measured dLLM inference on commodity CPUs.

This paper makes six contributions:

\begin{enumerate}
\item \textbf{First empirical demonstration} of the memory-compute regime inversion on
  CPU: we show that dLLMs achieve near-linear thread scaling ($7.5\times$ at 12 cores)
  while AR inference saturates at 4 cores ($2.4\times$).

\item \textbf{\texttt{diffuse-cpp}}: the first C++ inference engine for dLLMs, built on
  GGML~\cite{gerganov2023}, supporting two model architectures: LLaDA-8B
  (Llama backbone, MHA) and Dream-7B (Qwen2.5 backbone, GQA). Both use
  F16, Q8\_0, and Q4\_K\_M quantization.

\item \textbf{Entropy-exit adaptive scheduling}: an entropy-based early termination
  mechanism that reduces denoising steps from 16 to 2--4 on easy prompts, providing up
  to $6\times$ speedup over the F16 baseline with no quality degradation.

\item \textbf{Inter-step KV cache}: a caching mechanism that reuses Key and Value tensors
  between denoising steps, reducing per-step cost by only recomputing active positions.
  Achieves $1.6\times$ average speedup with an implicit regularization effect that
  improves output quality on 3 of 8 benchmarked prompts.

\item \textbf{Multi-model comparative benchmarking}: systematic comparison of two
  dLLM architectures (MHA vs GQA) on identical hardware, revealing complementary
  strengths---Dream-7B excels at math and code while LLaDA-8B excels at translation.

\item \textbf{Systematic benchmarking}: 288 configurations with 3 repetitions each,
  plus 8 diverse real prompts across 2 models, totaling over 8 hours of continuous
  CPU measurement.
\end{enumerate}


\section{Background}

\subsection{Autoregressive Inference}

AR models generate tokens sequentially: at each step $t$, the model computes
$\mathbf{y}_t = f(\mathbf{x}_{1:t})$ by multiplying the input embedding against the
weight matrices. For a single new token, this reduces to a matrix-vector product
$\mathbf{W} \cdot \mathbf{h}$ where $\mathbf{W} \in \mathbb{R}^{N \times K}$ and
$\mathbf{h} \in \mathbb{R}^{K}$. The arithmetic intensity (AI) is:
\begin{equation}
\text{AI}_{\text{AR}} = \frac{2NK}{(N+K) \cdot b} \approx \frac{2}{b} \quad \text{FLOP/byte}
\end{equation}
where $b$ is the bytes per parameter ($b=2$ for FP16, $b \approx 0.5$ for Q4\_K\_M).
For Q4\_K\_M, $\text{AI}_{\text{AR}} \approx 4$~FLOP/byte. KV caching eliminates
redundant computation but does not change the fundamental memory access pattern.

\subsection{Diffusion Inference}

Diffusion LLMs start with a fully masked sequence of $B$ tokens and refine it over $T$
steps. Each step performs a full forward pass on the entire sequence, producing
matrix-matrix products $\mathbf{W} \cdot \mathbf{H}$ where
$\mathbf{H} \in \mathbb{R}^{K \times B}$. The arithmetic intensity is:
\begin{equation}
\text{AI}_{\text{diff}} = \frac{2NKB}{(NK + KB) \cdot b} \approx \frac{2B}{b} \quad \text{FLOP/byte}
\end{equation}
For $B=64$ and Q4\_K\_M: $\text{AI}_{\text{diff}} \approx 256$~FLOP/byte; for $B=256$:
$\text{AI}_{\text{diff}} \approx 1024$~FLOP/byte---two to three orders of magnitude higher
than AR. Crucially, larger $B$ increases AI without adding memory traffic, so longer
generation buffers make diffusion \emph{more} compute-efficient per token (Section~\ref{sec:seqlen}).
There is no KV cache; each step recomputes attention over the full sequence with
bidirectional (non-causal) masking.

\subsection{Prior Work}

LLaDA~\cite{nie2025} introduced masked discrete diffusion for language modeling with a
Llama-style backbone and Multi-Head Attention (MHA). Dream~\cite{ye2025} adopted a
Qwen2.5 backbone with Grouped Query Attention (GQA), autoregressive logit shift for
causal bias, and additional remasking strategies. SEDD~\cite{lou2024} proposed
score-based discrete diffusion. MDLM~\cite{sahoo2024} simplified the training objective.
Prophet~\cite{li2025} introduced speculative decoding for discrete diffusion, using a
confidence-gap criterion for early stopping. EAGS~\cite{koh2024} and
Fast-dLLM~\cite{wu2025} explored entropy-adaptive scheduling. All prior work
benchmarked on GPU.

For CPU inference, \texttt{llama.cpp}~\cite{gerganov2023b} is the state of the art for
AR models, achieving near-theoretical throughput through GGML's optimized SIMD kernels
and quantization formats.


\section{The Regime Inversion}

The central insight of this paper is that the relative performance of AR and diffusion
inference \emph{inverts} depending on the hardware's memory-bandwidth-to-compute ratio.
We formalize this using the roofline model~\cite{williams2009}.

\subsection{The Matching Argument}

A processor achieves peak performance when the workload's arithmetic intensity matches
the machine's operational intensity (peak FLOPS / peak bandwidth). Table~\ref{tab:matching}
shows the matching:

\begin{table}[h]
\centering
\caption{Memory-compute matching across hardware regimes.}
\label{tab:matching}
\begin{tabular}{lcc}
\toprule
 & \textbf{GPU} (high BW) & \textbf{CPU} (low BW) \\
\midrule
\textbf{AR} (BW-hungry) & Match $\to$ fast & Mismatch $\to$ slow \\
\textbf{Diffusion} (compute-hungry) & Redundant & Match $\to$ scales \\
\bottomrule
\end{tabular}
\end{table}

On GPUs with $\sim$1~TB/s memory bandwidth, AR's low arithmetic intensity is tolerable:
the hardware delivers data fast enough. Diffusion's high compute demand, while achievable,
offers no advantage over highly optimized AR inference with KV caching.

On CPUs with $\sim$50--80~GB/s bandwidth, the situation reverses. AR inference hits a
hard ceiling: with a 5.1~GB model (Q4\_K\_M) and 75~GB/s bandwidth, the theoretical
maximum is $75/5.1 \approx 14.7$~tokens/s, and practical throughput is 8.51~tok/s
(58\% of theoretical). Adding threads provides diminishing returns because all cores
share the memory bus.

Diffusion inference, in contrast, is compute-bound on CPU. Each forward pass performs
$\mathcal{O}(NKB)$ operations with $\mathcal{O}(NK)$ data movement, and matrix-matrix
multiplication parallelizes efficiently across cores via BLAS routines. The throughput
scales with available compute, not bandwidth.

\subsection{Empirical Evidence: Thread Scaling}

Figure~\ref{fig:scaling} shows the critical empirical evidence. We benchmark both
\texttt{diffuse-cpp} (LLaDA-8B Q4\_K\_M, $T$=16 steps) and \texttt{llama.cpp}
(Llama-3-8B Q4\_K\_M) on identical hardware (AMD EPYC 4465P, 12 physical cores).

% Figure 1: Thread scaling comparison
% \begin{figure}[h]
% \centering
% \includegraphics[width=\textwidth]{fig1.pdf}
% \caption{Thread scaling comparison. (a) Absolute throughput: \texttt{llama.cpp} is
% faster at low core counts but saturates. (b) Normalized scaling: diffusion achieves
% $7.4\times$ at 12 threads versus $2.4\times$ for AR. Hyperthreading (24 threads)
% degrades both, confirming the physical-core bottleneck.}
% \label{fig:scaling}
% \end{figure}

The AR curve saturates at 4 threads (scaling from $1.0\times$ to $2.5\times$, then
stalling). The diffusion curve continues to $7.4\times$ at 12 threads, with degradation
only at 24 threads due to hyperthreading contention. This demonstrates compute-bound
scaling for diffusion: each additional core contributes proportionally to throughput,
as predicted by the roofline analysis.

\subsection{The Crossover}

In absolute terms, \texttt{llama.cpp} remains faster for raw token generation at 12
cores (8.51~tok/s for 16-step diffusion). However, this comparison is misleading for
two reasons. First, entropy-exit scheduling (Section~\ref{sec:entropy}) reduces effective
steps to 2--4 on favorable prompts, pushing diffusion throughput to 15--28~tok/s and
surpassing AR. Second, the scaling trends diverge: as core counts increase (modern
server CPUs have 64--128 cores), diffusion throughput grows near-linearly while AR
remains bandwidth-limited. The crossover for unconditional throughput likely occurs
between 16 and 32 cores on current hardware.


\section{\texttt{diffuse-cpp}: Architecture}

\texttt{diffuse-cpp} is a C++ inference engine for masked diffusion language models,
built on the GGML tensor library~\cite{gerganov2023}---the same foundation as
\texttt{llama.cpp}. This section describes its design.

\textbf{Supported architectures.} \texttt{diffuse-cpp} supports two model families:

\begin{itemize}
\item \textbf{LLaDA-8B}~\cite{nie2025}: Llama backbone, 32 layers, 4096 hidden dim,
  Multi-Head Attention (32 query / 32 KV heads), SwiGLU FFN (12288 intermediate),
  vocabulary 126{,}464, RoPE $\theta$=500{,}000.
\item \textbf{Dream-7B}~\cite{ye2025}: Qwen2.5 backbone, 28 layers, 3584 hidden dim,
  Grouped Query Attention (28 query / 4 KV heads), SwiGLU FFN (18944 intermediate),
  vocabulary 152{,}064, RoPE $\theta$=1{,}000{,}000. Dream adds QKV biases (kept at F32
  in all quantizations), an autoregressive logit shift for causal bias, and additional
  remasking strategies (\texttt{maskgit\_plus}, \texttt{topk\_margin}).
\end{itemize}

Both models use bidirectional (non-causal) attention. The shared forward pass differs
only in GQA handling: for Dream, each KV head is expanded to serve $n_\text{rep}=7$
query heads via grouped repetition before attention computation.

\textbf{Denoising loop.} Generation proceeds as follows: (1) initialize a sequence of
$B$ mask tokens; (2) perform a forward pass to obtain logits for all positions; (3) for
each masked position, sample from the logit distribution; (4) unmask positions based on
the scheduling policy; (5) repeat until all positions are unmasked or the step budget is
exhausted.

\textbf{Model format.} Models are stored in GGUF format, converted from HuggingFace
SafeTensors via dedicated Python scripts (\texttt{tools/convert-llada.py} and
\texttt{tools/convert-dream.py}). LLaDA produces 226 tensors; Dream produces 339
tensors (255 weights + 84 QKV biases).

\textbf{Quantization.} \texttt{diffuse-cpp} supports three quantization levels via
\texttt{diffuse-quantize}: F16, Q8\_0, and Q4\_K\_M. For LLaDA: 14.9/8.4/5.1~GB;
for Dream: 15.2/8.6/5.3~GB. Quantization is weight-only; activations remain in FP32.
QKV biases (Dream only) are preserved at F32 in all quantization levels.

\textbf{Scheduling policies.} Two scheduling policies are implemented.
\texttt{low\_confidence} is the default baseline: at each step, it unmasks a fixed
fraction of the most confident positions (lowest logit entropy).
\texttt{entropy\_exit} is our adaptive scheduler, described in
Section~\ref{sec:entropy}.

\subsection{Inter-Step KV Cache}
\label{sec:cache}

Standard diffusion inference recomputes the full attention for all $N$ positions at every
denoising step, even though most positions do not change between consecutive steps. Once
a position is unmasked with high confidence, its token value---and therefore its Key and
Value vectors---remains stable. We exploit this by caching the post-RoPE K and V tensors
from each transformer layer between steps.

\textbf{Mechanism.} After each full forward pass, we store the K and V tensors for all
$L$ layers and all $N$ positions. On subsequent steps, we partition positions into two sets:

\begin{itemize}
\item \textbf{Active set} $\mathcal{A}$: positions that are still masked or were recently
  unmasked. These require fresh K,V computation because their hidden states may have changed.
\item \textbf{Cached set} $\mathcal{C} = \{1,\ldots,N\} \setminus \mathcal{A}$: positions
  whose tokens have not changed. Their K,V are reused from the cache.
\end{itemize}

For each layer, the attention computation becomes:
\[
\text{Attn}(\mathbf{Q}_\mathcal{A}, [\mathbf{K}_\mathcal{C}^{\text{cached}} \| \mathbf{K}_\mathcal{A}^{\text{fresh}}], [\mathbf{V}_\mathcal{C}^{\text{cached}} \| \mathbf{V}_\mathcal{A}^{\text{fresh}}])
\]
where $\|$ denotes concatenation along the sequence dimension.

\textbf{Complexity.} The per-step cost reduces from $\mathcal{O}(N^2)$ per layer
(full recomputation) to $\mathcal{O}(|\mathcal{A}| \cdot N)$ per layer (only active
queries, full key-value context). Since $|\mathcal{A}|$ shrinks monotonically as tokens
converge, later steps become progressively cheaper.

\textbf{Approximation.} The cache introduces an approximation: the K,V for cached
positions are computed from \emph{stale} hidden states. In a bidirectional transformer,
changing any token affects all positions through attention interactions. However, we
observe empirically that this approximation does not degrade output quality. In fact,
it acts as implicit regularization: reusing stable K,V from earlier steps prevents
attention from amplifying accumulated errors during the denoising process. On 3 of 8
benchmarked prompts, cached output quality is measurably better than uncached
(Section~\ref{sec:cache_results}).

\textbf{Configuration.} The cache is enabled by default. The flag \texttt{--no-cache}
disables it. The flag \texttt{--cache-keep-active N} extends the active set to include
positions that changed within the last $N$ steps, trading speed for fidelity.


\section{Entropy-Exit Adaptive Scheduling}
\label{sec:entropy}

Standard diffusion schedulers use a fixed number of denoising steps $T$ regardless of
prompt difficulty. This is suboptimal: a factual question (``What is the capital of
France?'') requires far less refinement than creative writing (``Write a poem about the
ocean''). We introduce \texttt{entropy\_exit}, an adaptive scheduler that allows the
model to decide when to stop.

\subsection{Algorithm}

After each denoising step, \texttt{entropy\_exit} computes the Shannon
entropy~\cite{shannon1948} of the logit distribution at each remaining masked position:
\begin{equation}
H_i = -\sum_{v=1}^{V} p_{i,v} \log p_{i,v}
\end{equation}
where $p_{i,v} = \text{softmax}(\mathbf{z}_i)_v$ and $V$ is the vocabulary size.
Positions with $H_i < \tau$ (default $\tau = 1.5$~nats) are immediately
unmasked---the model is sufficiently confident. If all $B$ positions are unmasked,
generation terminates.

The computational overhead is negligible: entropy computation is $\mathcal{O}(BV)$,
while the forward pass is $\mathcal{O}(BNK)$ with $N, K \gg V$ for the dominant
attention and MLP layers.

\subsection{Prior Art}

Prophet~\cite{li2025} uses a confidence-gap criterion (difference between top-2 logits)
for speculative decoding of diffusion models. EAGS~\cite{koh2024} and
Fast-dLLM~\cite{wu2025} explore entropy-based scheduling. Our contribution is not the
scheduling concept itself, but its implementation in a CPU-optimized engine where the
memory-compute regime inversion makes adaptive scheduling particularly impactful:
reducing steps has a direct, near-linear effect on throughput in the compute-bound regime.

\subsection{Empirical Behavior}

% Figure 2: Entropy-exit step allocation
% \begin{figure}[h]
% \centering
% \includegraphics[width=\textwidth]{fig2.pdf}
% \caption{Entropy-exit step allocation across 8 real prompts. Easy prompts (factual,
% arithmetic, simple translation) converge in 2--4 steps. Hard prompts (creative,
% explanatory) use all configured steps. Numbers above bars show throughput in tok/s.}
% \label{fig:entropy}
% \end{figure}

Easy prompts (factual QA, arithmetic, simple translation) converge in 2--4 steps,
achieving 15--28~tok/s with inter-step cache. Medium-difficulty prompts (code generation,
complex translation) use 7--15 steps. Hard prompts (creative writing, open-ended
explanation) use the full 15--16 steps with no speedup but also no degradation relative
to the baseline.

In a separate PyTorch-based evaluation across 42 prompts in 12 categories, 55\% of
prompts achieved $>$20\% speedup. The best categories were instruction-following
($3.86\times$), classification ($2.84\times$), and translation ($2.71\times$). No
category showed degradation.


\section{Experimental Results}

\subsection{Setup}

All benchmarks were conducted on a dedicated AMD EPYC 4465P (12 cores, 24 threads),
125~GB DDR5 RAM, running Ubuntu 24.04 with GCC 13.3. No competing processes ran during
measurement. Each configuration was run 3 times with 1 warmup, generating 64 tokens from
a 32-token prompt unless otherwise noted. Sequence-length scaling experiments
(Section~\ref{sec:seqlen}) used 256 generated tokens. GGML compute buffers used a
$1.5\times$ safety multiplier per upstream recommendation.

\textbf{Methodology note.} Initial benchmarks used a $1.2\times$ buffer multiplier,
yielding approximately 30\% higher throughput (e.g., 13.59~tok/s peak). Upon discovering
this was below the GGML-recommended minimum, all 288 configurations were re-run with
$1.5\times$ buffers. This paper reports only the corrected results.

\subsection{Quantization Performance}

Table~\ref{tab:quant} shows the effect of quantization at a fixed configuration
($T$=16 steps, 12 threads).

\begin{table}[h]
\centering
\caption{Performance by quantization level (steps=16, threads=12, LLaDA-8B).}
\label{tab:quant}
\begin{tabular}{llrrr}
\toprule
\textbf{Model} & \textbf{Size} & \textbf{low\_confidence} & \textbf{Per-step} & \textbf{vs F16} \\
\midrule
F16 & 14.9\,GB & 1.64\,tok/s & 2433\,ms & $1.00\times$ \\
Q8\_0 & 8.4\,GB & 1.84\,tok/s & 2182\,ms & $1.12\times$ \\
Q4\_K\_M & 5.1\,GB & 2.52\,tok/s & 1586\,ms & $1.54\times$ \\
\bottomrule
\end{tabular}
\end{table}

The $1.54\times$ speedup from F16 to Q4\_K\_M is lower than the $2.9\times$ size
reduction would suggest. This is consistent with compute-bound behavior: reducing data
size helps (fewer bytes to dequantize), but the bottleneck is the matrix multiplication
itself, not the memory transfer.

\subsection{Thread Scaling}

Table~\ref{tab:threads} shows thread scaling for Q4\_K\_M at 16 steps.

\begin{table}[h]
\centering
\caption{Thread scaling (Q4\_K\_M, steps=16, \texttt{low\_confidence}).}
\label{tab:threads}
\begin{tabular}{rrr}
\toprule
\textbf{Threads} & \textbf{tok/s} & \textbf{Scaling} \\
\midrule
1 & 0.34 & $1.0\times$ \\
4 & 1.18 & $3.5\times$ \\
12 & 2.52 & $7.4\times$ \\
24 & 2.21 & $6.5\times$ \\
\bottomrule
\end{tabular}
\end{table}

Scaling is near-linear up to the physical core count (12), confirming compute-bound
behavior. Hyperthreading (24 threads) causes a 12\% regression due to shared execution
units. Optimal thread count equals the number of physical cores.

\subsection{Real-Prompt Performance}

Table~\ref{tab:real} shows end-to-end \texttt{entropy\_exit} performance on 8 diverse
prompts using Q4\_K\_M, $T$=16, 12 threads, at both $B$=64 and $B$=256, with and
without inter-step cache.

\begin{table}[h]
\centering
\caption{Real-prompt \texttt{entropy\_exit} performance (Q4\_K\_M, steps=16, threads=12).}
\label{tab:real}
\begin{tabular}{lrrrrrr}
\toprule
 & \multicolumn{2}{c}{$B$=64} & \multicolumn{3}{c}{$B$=256} \\
\cmidrule(lr){2-3} \cmidrule(lr){4-6}
\textbf{Prompt} & \textbf{tok/s} & \textbf{Steps} & \textbf{No-Cache} & \textbf{Cache} & \textbf{Steps} \\
\midrule
Capital of France? & 9.22 & 4 & 17.5 & \textbf{24.4} & 3 \\
Translate to French & 10.23 & 3 & 25.9 & \textbf{27.7} & 2 \\
$15 \times 23 = ?$ & 11.49 & 3 & 12.8 & \textbf{15.7} & 4 \\
Translate to Spanish & 4.59 & 8 & 7.6 & \textbf{22.9} & 7 \\
Python \texttt{is\_prime} & 2.53 & 15 & 3.2 & \textbf{4.9} & 16 \\
Poem about ocean & 2.33 & 17 & 3.2 & \textbf{5.3} & 16 \\
Why sky is blue & 2.21 & 17 & 3.3 & \textbf{12.0} & 16 \\
List the planets & 2.33 & 17 & 3.3 & \textbf{9.4} & 15 \\
\bottomrule
\end{tabular}
\end{table}

The pattern at $B$=64 is clear: easy prompts (factual, arithmetic, simple translation)
converge in 2--4 steps with $4$--$6\times$ speedup over \texttt{low\_confidence}. Hard
prompts use full steps with no speedup but no degradation.

At $B$=256, all prompts benefit from the larger generation buffer, with throughput gains
of 23--113\% on 7 of 8 prompts. With inter-step cache enabled (default), throughput
increases by an additional $1.59\times$ on average, with the highest gains on prompts
that require many steps (up to $3.66\times$ on 16-step prompts).

\subsection{Comparison with \texttt{llama.cpp}}

Table~\ref{tab:comparison} compares \texttt{diffuse-cpp} with \texttt{llama.cpp} on the
same hardware. We emphasize that this compares different models (LLaDA-8B vs Llama-3-8B)
and different paradigms, not identical systems.

\begin{table}[h]
\centering
\caption{Comparison with \texttt{llama.cpp} (both Q4\_K\_M, 12 threads, same hardware).}
\label{tab:comparison}
\begin{tabular}{lrrr}
\toprule
\textbf{Scenario} & \textbf{diffuse-cpp} ($B$=256, cache) & \textbf{llama.cpp} \\
\midrule
Easy prompts (EE) & \textbf{15--28}\,tok/s & 8.51\,tok/s \\
Medium prompts (EE) & 5--23\,tok/s & 8.51\,tok/s \\
Hard prompts (EE) & 5--12\,tok/s & 8.51\,tok/s \\
Thread scaling ($1 \to 12$) & $7.4\times$ & $2.4\times$ \\
\bottomrule
\end{tabular}
\end{table}

With $B$=256 and inter-step cache, \texttt{diffuse-cpp} outperforms \texttt{llama.cpp}
on 6 of 8 prompts (up to $3.3\times$ on French translation). Two prompts (code generation,
creative writing) remain slower due to requiring all 16 denoising steps. The thread
scaling trends further favor diffusion on many-core machines.

% Figure 3: Combined optimization
% \begin{figure}[h]
% \centering
% \includegraphics[width=\textwidth]{fig3.pdf}
% \caption{Combined optimization: from F16 baseline (1.64~tok/s) through quantization
% (Q4\_K\_M, $1.54\times$), entropy-exit on easy prompts at $B$=64 ($6.2\times$),
% $B$=256 generation ($13.3\times$), and inter-step cache ($16.9\times$ peak,
% 27.7~tok/s). The dashed line shows \texttt{llama.cpp} for reference.}
% \label{fig:combined}
% \end{figure}

\subsection{Sequence Length Scaling}
\label{sec:seqlen}

Table~\ref{tab:seqlen} shows the effect of increasing the generation buffer from 64 to
256 tokens (Q4\_K\_M, $T$=16 steps, 12 threads, real chat-formatted prompt).

\begin{table}[h]
\centering
\caption{Sequence length scaling (Q4\_K\_M, steps=16, threads=12, 32-token prompt).}
\label{tab:seqlen}
\begin{tabular}{lrrrrr}
\toprule
\textbf{Scheduler} & $B$=64 & $B$=256 & \textbf{Time ratio} & \textbf{ms/token} & \textbf{Improvement} \\
\midrule
\texttt{low\_confidence} & 2.58\,tok/s & 3.18\,tok/s & $3.24\times$ & $388 \to 314$ & $-19\%$ \\
\texttt{entropy\_exit} & 13.4\,tok/s & 17.1\,tok/s & $3.15\times$ & $74.5 \to 58.6$ & $-21\%$ \\
\bottomrule
\end{tabular}
\end{table}

Generating $4\times$ more tokens increases total wall-clock time by only $3.2\times$,
yielding a 20\% reduction in per-token cost. This sub-linear scaling confirms the
theoretical prediction from Equation~2: larger $B$ increases arithmetic intensity,
amortizing fixed overhead (weight loading, softmax over vocabulary, memory allocation)
across more useful tokens.

The practical implication is significant: \texttt{entropy\_exit} with $B$=256 reaches
\textbf{15--28~tok/s} on real prompts (Table~\ref{tab:real})---up to $3.3\times$ the
throughput of \texttt{llama.cpp} (8.51~tok/s) on the same hardware. This is the
configuration where diffusion's matrix-matrix advantage becomes most pronounced, and it
corresponds to the sequence lengths users actually need in production (responses of
200+ tokens are typical for chatbot, code generation, and translation tasks).

\subsection{Inter-Step Cache Results}
\label{sec:cache_results}

Table~\ref{tab:cache} shows the inter-step cache performance across all 8 real prompts
(Q4\_K\_M, \texttt{entropy\_exit}, $B$=256, $T$=16, 12 threads, seed=42).

\begin{table}[h]
\centering
\caption{Inter-step cache performance (Q4\_K\_M, entropy\_exit, $B$=256, steps=16, threads=12).}
\label{tab:cache}
\begin{tabular}{lrrrr}
\toprule
\textbf{Prompt} & \textbf{No-Cache} & \textbf{Cache} & \textbf{Speedup} & \textbf{Steps} \\
\midrule
Capital of France? & 17.5\,tok/s & \textbf{24.4}\,tok/s & $1.39\times$ & 3 \\
Translate to French & 25.9 & \textbf{27.7} & $1.07\times$ & 2 \\
Python \texttt{is\_prime} & 3.2 & \textbf{4.9} & $1.52\times$ & 16 \\
Poem about ocean & 3.2 & \textbf{5.3} & $1.63\times$ & 16 \\
Why sky is blue & 3.3 & \textbf{12.0} & $3.66\times$ & 16 \\
List the planets & 3.3 & \textbf{9.4} & $2.84\times$ & 15 \\
$15 \times 23 = ?$ & 12.8 & \textbf{15.7} & $1.22\times$ & 4 \\
Translate to Spanish & 7.6 & \textbf{22.9} & $3.03\times$ & 7 \\
\midrule
\textbf{Average} & \textbf{9.6} & \textbf{15.3} & $\mathbf{1.59\times}$ & \\
\bottomrule
\end{tabular}
\end{table}

The cache provides $1.59\times$ average speedup with a clear pattern: prompts requiring
more denoising steps benefit more. Two-step prompts gain only $1.07\times$ (the active
set is nearly the full sequence), while 16-step prompts gain up to $3.66\times$ (most
positions are cached after early steps).

\textbf{Regularization effect.} In 3 of 8 prompts (sky explanation, planet listing,
Spanish translation), the cached output is qualitatively better than the uncached
output. We attribute this to implicit regularization: reusing stable K,V from previous
steps prevents the bidirectional attention mechanism from amplifying accumulated errors
during the iterative denoising process. This phenomenon warrants further investigation
but suggests that the cache approximation may be beneficial rather than merely tolerable.

\textbf{Fidelity tuning.} The \texttt{--cache-keep-active N} parameter extends the
active set to include positions that changed within the last $N$ steps. At $N$=5, token
fidelity (exact match with uncached output) reaches 59.4\% with $1.64\times$ speedup---a
useful trade-off when output reproducibility matters.

\subsection{Multi-Model Comparison: Dream-7B vs LLaDA-8B}
\label{sec:multimodel}

Table~\ref{tab:multimodel} compares the two supported architectures on all 8 prompts
under identical conditions (Q4\_K\_M, \texttt{entropy\_exit} + inter-step cache, 12
threads, seed=42). Dream generates $B$=64 tokens; LLaDA generates $B$=256.

\begin{table}[h]
\centering
\caption{Dream-7B vs LLaDA-8B (Q4\_K\_M, entropy\_exit + cache, 12 threads). Bold indicates the faster model per prompt.}
\label{tab:multimodel}
\begin{tabular}{lrrrrrr}
\toprule
 & \multicolumn{2}{c}{\textbf{Dream-7B}} & \multicolumn{2}{c}{\textbf{LLaDA-8B}} & \multicolumn{2}{c}{\textbf{vs llama.cpp}} \\
\cmidrule(lr){2-3} \cmidrule(lr){4-5} \cmidrule(lr){6-7}
\textbf{Prompt} & \textbf{tok/s} & \textbf{Steps} & \textbf{tok/s} & \textbf{Steps} & \textbf{Dream} & \textbf{LLaDA} \\
\midrule
Capital of France?   & \textbf{21.6} & 2  & 22.4 & 3  & $2.5\times$ & $2.6\times$ \\
Translate to French  & 14.3 & 6  & \textbf{25.7} & 2  & $1.7\times$ & $3.0\times$ \\
$15 \times 23 = ?$   & \textbf{21.6} & 2  & 6.0  & 16 & $2.5\times$ & $0.7\times$ \\
Translate to Spanish & 13.2 & 10 & \textbf{23.3} & 5  & $1.6\times$ & $2.7\times$ \\
Python \texttt{is\_prime} & \textbf{8.2} & 7 & 4.5 & 15 & $1.0\times$ & $0.5\times$ \\
Why sky is blue      & 4.9  & 16 & \textbf{5.0}  & 16 & $0.6\times$ & $0.6\times$ \\
List the planets     & 4.9  & 16 & \textbf{9.5}  & 16 & $0.6\times$ & $1.1\times$ \\
Poem about ocean     & 4.5  & 16 & \textbf{5.0}  & 16 & $0.5\times$ & $0.6\times$ \\
\midrule
\textbf{Average}     & \textbf{11.6} & & \textbf{12.7} & & $1.4\times$ & $1.5\times$ \\
\bottomrule
\end{tabular}
\end{table}

\textbf{Complementary strengths.} The two models exhibit complementary performance
profiles. Dream-7B excels at \emph{math and code}: it correctly solves $15 \times 23 = 345$
in 2 steps (21.6~tok/s), where LLaDA requires all 16 steps and produces an incorrect
answer. Dream also generates a complete, correct \texttt{is\_prime()} function in 7
steps (8.2~tok/s) versus LLaDA's 15 steps (4.5~tok/s). LLaDA-8B excels at
\emph{translation}: French translation converges in 2 steps at 25.7~tok/s ($3.0\times$
llama.cpp), while Dream requires 6 steps (14.3~tok/s). Both models struggle with
creative writing (4.5--5.0~tok/s), requiring all 16 steps.

\textbf{Architecture effects.} Dream's GQA (28/4 heads) reduces KV computation by
$7\times$ compared to LLaDA's MHA (32/32 heads), but this advantage is partially offset
by Dream's smaller generation buffer ($B$=64 vs $B$=256). The inter-step cache benefits
both architectures; for Dream, cached K,V are stored post-GQA-expansion (at full
$n_\text{head}$ dimensionality) to avoid recomputing the grouped repetition on each step.

\textbf{Practical recommendation.} Users should select the model based on their primary
use case: Dream for math, code, and factual queries; LLaDA for translation and
conversational tasks. Both outperform \texttt{llama.cpp} on 5--6 of 8 prompts.


\section{Limitations}

We are transparent about the limitations of this work.

\textbf{Model quality.} Both LLaDA-8B and Dream-7B produce lower-quality outputs than
their AR counterparts (Llama-3-8B, Qwen2.5-7B). LLaDA shows occasional repetitions and
incomplete code; Dream produces higher-quality math and code but sometimes generates
hallucinated content. The throughput comparison is between paradigms, not
quality-equivalent models.

\textbf{Different models.} We compare dLLMs (LLaDA-8B, Dream-7B) against AR models
(Llama-3-8B). These share similar parameter counts and backbone architectures but differ
in training procedure, data, and generation mechanism. The comparison demonstrates the
\emph{computational} regime inversion, not a head-to-head model quality contest.

\textbf{Generation buffer asymmetry.} Dream benchmarks use $B$=64 tokens while LLaDA
uses $B$=256. This reflects practical constraints (Dream's Qwen2.5 tokenizer produces
shorter outputs) but means the per-token throughput comparison is not perfectly
controlled. Dream would likely achieve higher throughput at $B$=256.

\textbf{Prior art on scheduling.} Entropy-based scheduling is not novel.
Prophet~\cite{li2025}, EAGS~\cite{koh2024}, and Fast-dLLM~\cite{wu2025} proposed
related mechanisms. Our contribution is the implementation in a CPU-optimized engine
and the demonstration that adaptive scheduling is particularly impactful in the
compute-bound regime.

\textbf{Single hardware platform.} All benchmarks were conducted on a single AMD EPYC
system. Generalization to Intel, ARM, Apple Silicon, or consumer-grade CPUs requires
further validation. The regime inversion argument is architecture-independent, but the
crossover point varies with the bandwidth-to-compute ratio.

\textbf{Prompt coverage.} Real-prompt evaluation used 8 prompts. While the PyTorch
benchmark covered 42 prompts in 12 categories, comprehensive evaluation across thousands
of diverse prompts remains future work.

\textbf{Cache approximation.} The inter-step KV cache reuses K,V tensors computed from
stale hidden states. While empirically this does not degrade quality (and sometimes
improves it), the approximation error grows with the number of steps between cache
refresh. For very long generation sequences or models with stronger inter-position
coupling, the cache may require periodic full refreshes.


\section{Implications and Future Work}

\textbf{Democratizing inference.} If diffusion models perform better on commodity CPUs,
the implications are significant. A \$200 server CPU with 32--64 cores could provide
competitive LLM inference without any GPU. This eliminates the hardware supply chain
bottleneck and enables deployment in resource-constrained environments.

\textbf{AR-to-diffusion conversion.} Recent work shows that AR models can be converted
to diffusion models via continued pre-training with masking~\cite{nie2025}. We estimate
that LoRA fine-tuning for this conversion could cost as little as \$250 in compute,
creating a practical pipeline: take any AR model, fine-tune for diffusion, convert to
GGUF, and deploy on CPU with \texttt{diffuse-cpp}.

\textbf{Multi-hardware validation.} A community hackathon or Kaggle competition could
validate the regime inversion across diverse hardware: Intel Xeon, AMD Threadripper,
Apple M-series, ARM server chips. We predict that the advantage of diffusion will
increase with core count and decrease with memory bandwidth.

\textbf{Engine optimizations.} The inter-step KV cache (Section~\ref{sec:cache})
demonstrates that significant speedups are achievable through caching strategies.
Further optimizations remain unexplored: bidirectional Flash Attention~\cite{dao2022}
adapted for non-causal masks, delta attention (recomputing only changed hidden state
contributions), batched inference for serving workloads, generation buffers beyond 256
tokens, and support for additional architectures (SEDD, MDLM). Dream-7B's GQA
architecture also opens the possibility of a more memory-efficient inter-step cache
that stores K,V at the reduced KV-head dimensionality rather than post-expansion.

\textbf{Better diffusion models.} The regime inversion becomes practically relevant when
dLLM quality matches AR models. Investment in larger, better-trained diffusion models
(analogous to the Llama series) would enable the computational advantages demonstrated
here to translate into real-world deployment wins.


\section{Conclusion}

The field has assumed that diffusion language models are inherently slower than
autoregressive models for inference. This assumption is correct on GPUs, where massive
memory bandwidth compensates for AR's sequential memory access pattern. On commodity
CPUs, however, the computational regime inverts: AR inference hits a hard bandwidth
ceiling while diffusion inference scales with available compute.

We demonstrated this empirically with \texttt{diffuse-cpp}, the first C++ inference
engine for diffusion LLMs, supporting two architectures: LLaDA-8B (Llama/MHA) and
Dream-7B (Qwen2.5/GQA). On an AMD EPYC 12-core CPU, thread scaling reaches
$7.4\times$ for diffusion versus $2.4\times$ for AR. With Q4\_K\_M quantization,
entropy-exit scheduling, and inter-step KV cache, \texttt{diffuse-cpp} achieves
14--28~tok/s on easy prompts---up to $3.3\times$ faster than \texttt{llama.cpp}'s
8.51~tok/s on the same hardware. The two models exhibit complementary strengths:
Dream excels at math and code (21.6~tok/s), while LLaDA excels at translation
(27.7~tok/s). The inter-step cache provides $1.6\times$ average speedup with no quality
degradation.

The contribution is not a claim that diffusion is universally faster. It is the
identification and empirical validation of a hardware-dependent regime inversion, and the
open-source tools to exploit it. As CPU core counts increase and diffusion model quality
improves, this regime may become the dominant path to GPU-free LLM inference.

The code, models, benchmark data, and this paper are publicly available under open licenses.


\section*{Acknowledgments}

This work was inspired by the GGML library and llama.cpp~\cite{gerganov2023b}. The
entropy-exit scheduler builds on ideas from Prophet~\cite{li2025} and the
entropy-adaptive scheduling literature~\cite{koh2024,wu2025}. Dream support was
developed using the architecture specifications from~\cite{ye2025}. Figures~1--3 were
generated with Matplotlib.


\bibliographystyle{plainnat}

\begin{thebibliography}{12}

\bibitem[Dao et~al.(2022)]{dao2022}
Tri Dao, Dan Fu, Stefano Ermon, Atri Rudra, and Christopher R\'e.
\newblock FlashAttention: Fast and memory-efficient exact attention with
  IO-awareness.
\newblock In \emph{Advances in Neural Information Processing Systems}, 2022.

\bibitem[Gerganov(2023)]{gerganov2023}
Georgi Gerganov.
\newblock GGML: Tensor library for machine learning.
\newblock \url{https://github.com/ggerganov/ggml}, 2023.

\bibitem[Gerganov et~al.(2023)]{gerganov2023b}
Georgi Gerganov et~al.
\newblock llama.cpp.
\newblock \url{https://github.com/ggerganov/llama.cpp}, 2023.

\bibitem[Koh et~al.(2024)]{koh2024}
Hyukhun Koh, Minha Jhang, Dohyung Kim, Sangmook Lee, and Kyomin Jung.
\newblock PLM-based discrete diffusion language models with entropy-adaptive
  Gibbs sampling.
\newblock \emph{arXiv preprint arXiv:2411.06438}, 2024.

\bibitem[Li et~al.(2025)]{li2025}
Pengxiang Li, Yefan Zhou, Dilxat Muhtar, Lu Yin, Shilin Yan, Li Shen, Soroush
  Vosoughi, and Shiwei Liu.
\newblock Diffusion language models know the answer before decoding.
\newblock \emph{arXiv preprint arXiv:2508.19982}, 2025.

\bibitem[Lou et~al.(2024)]{lou2024}
Aaron Lou, Chenlin Meng, and Stefano Ermon.
\newblock Discrete diffusion modeling by estimating the ratios of the data
  distribution.
\newblock In \emph{International Conference on Machine Learning (ICML)}, 2024.

\bibitem[Nie et~al.(2025)]{nie2025}
Shen Nie, Fengqi Zhu, Chao You, Xiaojun Zhang, and Zhenguo Ou.
\newblock LLaDA: Large language diffusion with mAsking.
\newblock \emph{arXiv preprint arXiv:2502.09992}, 2025.

\bibitem[Sahoo et~al.(2024)]{sahoo2024}
Subham Sekhar Sahoo, Marianne Arriola, Yair Schiff, Aaron Gokaslan, Edgar
  Marroquin, Justin~T Chiu, Alexander Rush, and Volodymyr Kuleshov.
\newblock Simple and effective masked diffusion language models.
\newblock \emph{Advances in Neural Information Processing Systems}, 2024.

\bibitem[Shannon(1948)]{shannon1948}
Claude~E Shannon.
\newblock A mathematical theory of communication.
\newblock \emph{The Bell System Technical Journal}, 27(3):379--423, 1948.

\bibitem[Williams et~al.(2009)]{williams2009}
Samuel Williams, Andrew Waterman, and David Patterson.
\newblock Roofline: An insightful visual performance model for multicore
  architectures.
\newblock \emph{Communications of the ACM}, 52(4):65--76, 2009.

\bibitem[Wu et~al.(2025)]{wu2025}
Chengyue Wu, Hao Zhang, Shuchen Xue, Zhijian Liu, Shizhe Diao, Ligeng Zhu,
  Ping Luo, Song Han, and Enze Xie.
\newblock Fast-dLLM: Training-free acceleration of diffusion LLM by enabling KV
  cache and parallel decoding.
\newblock \emph{arXiv preprint arXiv:2505.22618}, 2025.

\bibitem[Ye et~al.(2025)]{ye2025}
Jiacheng Ye, Shansan Gong, Liheng Chen, Lin Zheng, Mukai Li, Thong Nguyen,
  Kaiyan Zhang, Shuicheng Yan, Silvio Savarese, Caiming Xiong, and Yingbo Zhou.
\newblock Dream 7B -- Scalable generation of structured and unstructured content
  from diffusion language modeling.
\newblock \emph{arXiv preprint arXiv:2504.16915}, 2025.

\end{thebibliography}


\appendix

\section{Regime Inversion Diagram}

% Figure 4: Regime inversion diagram
% \begin{figure}[h]
% \centering
% \includegraphics[width=0.8\textwidth]{fig4.pdf}
% \caption{The $2 \times 2$ matching matrix. Green cells indicate favorable
% hardware-workload matching; red cells indicate bottlenecks. On GPU, AR is
% well-matched; on CPU, diffusion is well-matched.}
% \label{fig:regime}
% \end{figure}

The $2 \times 2$ matching matrix illustrates the regime inversion:

\begin{table}[h]
\centering
\begin{tabular}{lcc}
\toprule
 & \textbf{GPU} (High Bandwidth) & \textbf{CPU} (Low Bandwidth) \\
\midrule
\textbf{Autoregressive} (BW-hungry) & \textsc{Match} & \textsc{Mismatch} \\
 & High BW feeds sequential reads & Low BW bottlenecks sequential reads \\
\midrule
\textbf{Diffusion} (Compute-hungry) & \textsc{Redundant} & \textsc{Match} \\
 & Excess BW wasted on parallel compute & Parallel compute fills CPU cores \\
\bottomrule
\end{tabular}
\end{table}


\end{document}
